
\section{“\trPoincare”对偶}
\label{sec:poincare-duality}

我们知道，$f$ 的\trIndex{}为 $k$ 的诸\trCriticalPoint{}恰是 $-f$ 的\trIndex{}为 $n-k$
的诸\trCriticalPoint（这里 $n$ 是\trManifold{}的维数）。此外，若 $X$ 是\trAdapted{}于
$f$ 的一个\trPseudogradient，那么 $-X$ 就是\trAdapted{}于 $-f$ 的一个\trPseudogradient。\footnote{译注：本节标题给“\trPoincare”加引号系原书所有——这里建立的是%
\trMorseComplex{}层面（借助函数 $-f$）的对偶性，而非直接援引经典的%
\trPoincareDuality{}定理。}

若 $a\in\operatorname{Crit}_k(f)$，我们令 $a^\star$ 表示同一个点，只是把它看成函数 $-f$
的\trCriticalPoint，于是
\[
\{a \mid a\in\operatorname{Crit}_k(f)\}\ \text{是 } C_k(f)\ \text{的一个基}
\]
而
\[
\{a^\star \mid a\in\operatorname{Crit}_k(f)\}\ \text{是 } C_{n-k}(-f)\ \text{的一个基。}
\]
\trVectorSpace{} $C_{n-k}(-f)$ \trIsomorphic{}于 $C_k(f)$，但更恰当的想法是把它看成后者的
\textbf{对偶}（dual，记号即由此而来）：
\[
C_{n-k}(-f)=C_k(f)^\star\text{。}
\]
于是，按定义，$\partial_X\colon C_k(f)\longrightarrow C_{k-1}(f)$ 的转置
\[
{}^t\partial_X\colon C_{n-k+1}(-f)\longrightarrow C_{n-k}(-f)
\]
就是\trComplex{} $(C_\star(-f),\partial_{-X})$ 的\trDifferential{} $\partial_{-X}$。
于是我们有下述结果。

\begin{proposition}[\trPoincareDuality]\label{thm:poincare-duality-mod2}\BookIndexDefinition{trPoincareDuality}
设 $V$ 是一个维数为 $n$ 的\trCompact{}\trManifold。若 $V$ 无边界，则
$HM_{n-k}(V;\ZZ/2)$ \trIsomorphic{}于 $HM_k(V;\ZZ/2)$。若 $V$ 有边界
$\partial V=\partial_+V\cup\partial_-V$，则 $HM_k(V,\partial_+V;\ZZ/2)$ \trIsomorphic{}于
$HM_{n-k}(V,\partial_-V;\ZZ/2)$。\noqed
\end{proposition}

例如（这是 $\partial_+V=\varnothing$ 的情形），
\[
HM_k(V;\ZZ/2)\ \text{\trIsomorphic{}于}\ HM_{n-k}(V,\partial V;\ZZ/2)\text{。}
\]

为了对整系数\trHomology{}陈述一个类似的结果，我们必须能够给 $f$ 的任意一个\trIndex
为 $k$ 的\trCriticalPoint{} $a$（\emph{连同 $W^s(a)$ 的一个\trOrientation}）配上 $-f$ 的一个%
\trIndex{}为 $n-k$ 的\trCriticalPoint{} $a^\star$（\emph{连同 $W^s(a^\star)=W^u(a)$ 的一个%
\trOrientation}）。为了简化陈述，我们因而取定\trManifold{} $V$ 的一个\trOrientation：
此时 $W^s(a)$ 的一个\trOrientation{}便给出 $W^u(a)$ 的一个\trOrientation（而不只是它的
一个\trCoorientation）。前面所用的论证依然有效，于是对于\trHomologyGroup{}的自由部分，
我们得到一个类似的结果。

\begin{proposition}[\trOrientation{}\trManifold{}的\trPoincareDuality]%
\label{thm:poincare-duality-oriented}
设 $V$ 是一个维数为 $n$ 的\trCompact{}\trOrientation{}\trManifold。若 $V$ 无边界，则%
\trAbelianGroup{} $HM_{n-k}(V;\ZZ)$ 与 $HM_k(V;\ZZ)$ 有相同的\trRank。若 $V$ 有边界
$\partial V=\partial_+V\cup\partial_-V$，则\trAbelianGroup{} $HM_k(V,\partial_+V;\ZZ)$
与 $HM_{n-k}(V,\partial_-V;\ZZ)$\linebreak 有相同的\trRank。\noqed
\end{proposition}

\noindent\textbf{注记。}$HM_k$ 与 $HM_{n-k}$ 的\trTorsion{}部分可以不同，正如我们在
$V=\CP^3(\RR)$ 的情形所注意到的那样：此时 $HM_1(\CP^3(\RR);\ZZ)=\ZZ/2$，而
$HM_2(\CP^3(\RR);\ZZ)=0$。\footnote{译注：原书此处印作 $H_2(\CP^3(\RR);\ZZ)$；按本
注记上下文所论均为\trMorseHomology{}的\trTorsion{}比较，疑应作
$HM_2(\CP^3(\RR);\ZZ)$；此处径改。——译者注}
至于 $HM_\star(\CP^n(\RR);\ZZ)$ 的计算，参见\cref{ex:ch04-25}。
